\documentclass[12pt]{article}
\usepackage[utf8]{inputenc}
\usepackage[T1]{fontenc}
\usepackage[margin=1in]{geometry}
\usepackage{setspace}
\usepackage{xcolor}
\usepackage[colorlinks=true, urlcolor=blue!60!black, linkcolor=blue!60!black]{hyperref}
\setlength{\parindent}{0em}
\setlength{\parskip}{0.45em}
\renewcommand{\baselinestretch}{1.0}
\begin{document}
\begin{center}
{\Large \textit{USAID v.\ Alliance for Open Society International, Inc.}, 570 U.S.\ 205 (2013)\par}
\vspace{0.3em}
{\normalsize Case Brief --- One-Page Edition\par}
\end{center}

\textbf{Memory Jogger.} Congress may define what it subsidizes, but it cannot make domestic grantees espouse the Government's viewpoint as their own outside the funded program (\href{https://supreme.justia.com/cases/federal/us/570/205/}{570 U.S. at 214--15, 221}).

\textbf{Facts.} The Leadership Act of 2003 funded NGOs fighting HIV/AIDS abroad but barred funds to any group without a policy explicitly opposing prostitution and sex trafficking. A separate, unchallenged provision already barred using the funds to promote prostitution. Domestic NGOs wished to stay neutral to keep the trust of host governments and sex workers.

\textbf{Procedural History.} The district court enjoined the Policy Requirement; the Second Circuit affirmed (2--1), splitting with the D.C. Circuit.

\textbf{Votes.} Affirmed, 6--2: Roberts for the Court; Scalia, with Thomas, dissented; Kagan took no part.

\textbf{Holding.} The Policy Requirement violates the First Amendment as applied to domestic organizations, because it regulates speech outside the program rather than defining the program.

\textbf{Reasoning.} Under \textit{Regan}, \textit{League of Women Voters}, and \textit{Rust}, the Government may limit what a funded project says but not what the grantee says on its own dime. Since funds already could not promote prostitution, the pledge went further, ``defining the recipient.'' Affiliates were no cure: a separate one cannot speak for the recipient, and a closely identified one forces ``evident hypocrisy.'' Scalia answered that the Government may choose allies and coerces no one who declines the grant. \emph{Later:} \textit{AOSI II} (2020) let the pledge apply to foreign affiliates.

\textbf{Examples.} \emph{Same-Side:} overdose-clinic grants conditioned on opposing drug decriminalization. \emph{Opposite-Side:} a clean-energy curriculum requirement confined to the funded program. \emph{Fence-Sitter:} anti-smoking grants denied to an NGO with tobacco-industry sponsors.

\textbf{Critique.} The Court admitted the line is ``hardly clear,'' and it left open whether the Government may simply prefer like-minded applicants.

\textbf{Key Quotation.} ``In the present context, the relevant distinction that has emerged from our cases is between conditions that define the limits of the government spending program---those that specify the activities Congress wants to subsidize---and conditions that seek to leverage funding to regulate speech outside the contours of the program itself.'' --- Chief Justice Roberts, 570 U.S. at 214--15.

\end{document}
